\documentclass[12pt]{article}
\newcommand{\bottomMargin}{2cm}

\newcommand{\header}{
	ДИЗАЙН И АНАЛИЗ НА АЛГОРИТМИ - \\ ПРАКТИКУМ\\
	\vspace{0.1cm}
	Летен семестър, 2026 г., трето контролно\\
	\vspace{0.1cm}
}
\newcommand{\tl}{$0.1$ сек.}
\newcommand{\ml}{$256$ MB}

\input{structure.tex}
\usepackage{newunicodechar}
\newunicodechar{ѝ}{{$\grave{\text{и}}$}}
\newunicodechar{Ѝ}{{$\grave{\text{И}}$}}

\begin{document}
	
\section{Задача K8. РАСПБЕРИ}

Котката Ксориндж иска да избяга от своя смъртен враг - чудовището Распбери. В момента тя е в капан, като се намира в квадратна таблица $N \times M$, изпълнена с храна (поставена от Распбери с цел да забави бягството на Ксориндж). Редовете на таблицата са номерирани с числата от $1$ до $N$, колоните са номерирани с числата от $1$ до $M$, а котката в началото се намира в клетка $(sx,sy)$. Храната в клетка $(i,j)$ носи $a_{i,j}$ удоволствие на Ксориндж, като е възможно някои храни да носят отрицателно удоволствие. Когато котката е в някоя клетка, тя винаги изяжда храната в нея.

Понеже чудовището тръгва от позиция, намираща се на северозапад спрямо Ксориндж, то котката се движи в посока югоизток. Това означава, че ако в момента тя е на клетка $(x,y)$, тя може да скочи извън таблицата, с което успява да избяга от капана, или да скочи в някоя от клетките в посока югоизток от $(x,y)$. По-точно, може да скочи в клетките с номера $(a,b) \neq (x,y)$, за които $x \le a \le N$ и $y \le b \le M$. Разбира се, като всяка котка, Ксориндж е ненаситна - тя така ще се пробва да избяга от таблицата, че да максимизира изпитаното удоволствие. Затова напишете програма \textbf{\texttt{raspberry}}, която за всяка възможна стартова клетка $(sx, sy)$ намира максималното удоволствие, което Ксориндж може да изпита, докато бяга от Распбери и се движи по описаните правила. За да провери вашите способности (и да не гледа ужасно много числа), котката иска от Вас само сумата на тези числа.

\subsection{Вход}
От първия ред на стандартния вход се въвеждат естествените числа $N$ и $M$ -- броят редове и колони на таблицата. От всеки от следващите $N$ реда се въвеждат по $M$ цели числа, които задават удоволствията на храната на съответния ред.

\subsection{Изход}
Изведете едно число - сумата от максималните удоволствия при бягството на Ксориндж за всяка една възможна стартова клетка $(sx, sy)$.

\subsection{Ограничения}
\vspace{0.1em}
\begin{itemize}[label=$\clubsuit$]
	\item $1 \leq N, M \leq 500$
	\item $-1000 \le a_{i,j} \le 1000$
\end{itemize}

\newpage

\subsection{Подзадачи}
\begin{table}[!ht]
	\begin{tblr}{|Q[c,m]|Q[c,m]|Q[c,m]|Q[c,m]|X[30,c,m]|}
		\hline
		\textbf{Подзадача} & \textbf{Точки} & $N$ & $M$ & \textbf{Други ограничения}\\
		\hline
		$1$ & $0$ & $-$ & $-$ & Примерният тест. \\ 
		\hline
		$2$ & $19$ & $=1$ & $\leq 500$ & За всяка стартова клетка, изпитаното максимално удоволствие включва преминаването през клетка $(N,M)$. \\ 
		\hline
		$3$ & $24$ & $\leq 5$ & $\leq 5$ & За всяка стартова клетка, изпитаното максимално удоволствие включва преминаването през клетка $(N,M)$. \\ 
		\hline
		$4$ & $45$ & $\leq 100$ & $\leq 100$ & За всяка стартова клетка, изпитаното максимално удоволствие включва преминаването през клетка $(N,M)$. \\ 
		\hline
		$5$ & $21$ & $\leq 500$ & $\leq 500$ & $a_{i,j} \ge 0$ \\ 
		\hline
		$6$ & $78$ & $\leq 500$ & $\leq 500$ & За всяка стартова клетка, изпитаното максимално удоволствие включва преминаването през клетка $(N,M)$. \\
		\hline 
		$7$ & $100$ & $\leq 500$ & $\leq 500$ & $-$ \\
		\hline
	\end{tblr}
    	\caption*{\textbf{Подзадача означава, че има тестове, носещи съответния брой точки, които спазват специфичните ограничения от подзадачата.}}
\end{table}
\FloatBarrier

\subsection{Пример}
\begin{table}[!ht]
	\begin{tblr}{|X[20,l]|X[8,l]|X[72,j]|}
		\hline
		\textbf{Вход} & \textbf{Изход} & \textbf{Обяснение на примера} \\
		\hline
		\texttt{4 3 \\
			1 2 3\\
			-2 -3 0\\
			-1 4 -2\\
			3 -1 -1}
		& 
		\texttt{25}
		& 
		{В следната таблица са показани максималните удоволствия, ако Ксориндж започне от всяка клетка:\\
		\begin{tabular}{|c|c|c|c|}
			\hline
			\texttt{7} & \texttt{6} & \texttt{3}\\
			\hline
			\texttt{2} & \texttt{1} & \texttt{0}\\
			\hline
			\texttt{3} & \texttt{4} & \texttt{-2}\\
			\hline
			\texttt{3} & \texttt{-1} & \texttt{-1}\\
			\hline
		\end{tabular}\\
		
		За да постигне максимално удоволствие, започвайки от клетка $(2,1)$, пътят на котката трябва да е: $(2,1) \to (3,2)$, след което тя скача извън таблицата и излиза от капана. Натрупаното удоволствие ще е $(-2)+4=2$.} \\
		\hline
	\end{tblr}
\end{table}
\FloatBarrier


	
\end{document}